\section{Linguaggi liberi da contesto}
I \textbf{linguaggi liberi da contesto} (Context-Free, CF) sono un insieme più generale di linguaggi che necessita di una specifica \textbf{grammatica}
per essere generato. Questa si definisce come una quadrupla $G = \left<V, T, P, S\right>$, dove: \begin{itemize}
    \item $V$: Insieme finito di simboli non terminali.
    \item $T$: Insieme finito di simboli terminali, tale che $V\cap T = \emptyset$.
    \item $P$: Insieme finito di produzioni, le quali hanno forma $A\to\alpha$, dove $A\in V$ è un simbolo non terminale e $\alpha\in (V\cup T)^*$.
        Quest'ultima è sostituibile a una qualunque occorrenza di $A$.
    \item $S\in V$: Simbolo iniziale.
\end{itemize}

\noindent Supponiamo per esempio $G = \left<\{E\}, \{0, 1, (, ), AND, OR, NOT\}, P, E \right>$, i cui elementi nelle seconde graffe rappresentano
l'alfabeto del linguaggio generato. Avremo le seguenti produzioni: \[ \begin{cases}
    E\to 0\\
    E\to 1\\
    E\to (E \text{ AND } E)\\
    E\to (E \text{ OR } E)\\
    E\to (\text{NOT } E)
\end{cases} \]

\noindent Scrittura che può essere contratta come: \[
    E\to 0|1|(E \text{ AND } E)|(E \text{ OR } E)|(\text{NOT } E)
\]

\noindent Naturalmente, le grammatiche servono a generare linguaggi con le seguenti relazioni: \[
    \xRightarrow{G}, \xRightarrow{G}_i, \xRightarrow{G}_* \subseteq (V\cup T)^*\times (V\cup T)^*
\]

\noindent Dove se $A\to \beta\in P$ e $\alpha, \gamma\in (V\cup T)^*$, allora $\alpha A \gamma \xRightarrow{G} \alpha \beta \gamma$. Diciamo qui che
$\alpha A \gamma$ deriva immediatamente $\alpha \beta \gamma$.\par
Inoltre, finché dentro $\alpha, \beta, \gamma$ sono presenti simboli non terminali, la generazione può continuare fino a ottenere una stringa in $T^*$.
Possiamo ora fornire la definizione di \textbf{linguaggio generato} dalla grammatica $G$: \begin{itemize}
    \item $\xRightarrow{G}_i$: Applicazione di $i$ passi di derivazione, ovvero di $i$ produzioni.
    \item $\xRightarrow{G}^*$: Esistenza di un numero di passi di derivazione non specificato.
    \item $\alpha\implies_i \beta$: Da $\alpha$ si deriva o genera $\beta$ in $i$ passi
    \item $\alpha\implies_* \beta$: Da $\alpha$ è possibile derivare o generare $\beta$.
\end{itemize}

\noindent Il linguaggio generato dalla grammatica $L(G) = \{x\in T^* | S\xRightarrow{G}_* x\}$ è l'insieme di tutte le stringhe in $T^*$ generate nella
grammatica, il quale si definisce \textbf{libero da cntesto} (CF) se esiste una grammatica altrettanto libera da contesto dotata di simboli terminali
$\Sigma$ capace di generarlo, facendo quindi valere $L = L(G)$.\par
Inoltre, due grammatiche $G_1, G_2$ si dicono equivalenti quando entrambe riescono a generare uno stesso linguaggio, dunque $L(G_1) = L(G_2)$. La
procedura di dimostrazione segue due passi: \begin{enumerate}
    \item $L\subseteq L(G)$: Conferma che tutte le stringhe di $L$ sono generabili in $G$. Ciò si fa per induzione sulla lunghezza di una stringa $x$:
    \[
        x\in L \implies S\implies_* x
    \]
    \item $L(G)\subseteq L$: Conferma che tutte le stringhe generabili in $G$ sono in $L$. Dimostrato per induzione sul numero di passi $i$: \[
        S\implies_i x \implies x\in L
    \]
\end{enumerate}

\noindent Le derivazioni generate dalle grammatiche sono visualizzabili appropriatamente grazie a strutture ad albero, dette \textbf{alberi di
derivazione} (parse tree), i quali, data una grammatica CF $G = \left<V, T, P, S\right>$, sono tali che: \begin{itemize}
    \item Ogni vertice ha un'etichetta in $V\cup T\cup \{\epsilon\}$.
    \item L'etichetta della radice è nei simboli non terminali $V$.
    \item Ogni vertice interno ha etichetta nei simboli non terminali $V$.
    \item Se un vertice $n$ ha etichetta in $A$ e $n_1, ..., n_k$ sono i suoi figli ordinati da sinistra a destra, etichettati con $X_1, ..., X_k$, allora $A\to X_1...X_k\in P$.
    \item Se un vertice ha etichetta $\epsilon$ allora è una foglia ed è l'unico figlio di suo padre.
\end{itemize}

\noindent Per esempio, supponiamo una produzione $E\to 0|1|(E \text{ OR } E)|(E \text{ AND } E)|(\text{NOT } E)$, avremo il seguente albero di
derivazione: \begin{figure}[h]
    \centering
    \begin{tikzpicture}[    
        level distance=1.5cm,
        level 1/.style={sibling distance=2cm},
        level 2/.style={sibling distance=1.3cm},
        level 3/.style={sibling distance=1.3cm},
        level 4/.style={sibling distance=1.3cm}
    ]

    \node {$E$}
        child {node {$($}}
        child {node {$E$}
            child {node {$($}}
            child {node {$E$}
                child {node {$0$}}
            }
            child {node {$OR$}}
            child {node {$E$}
                child {node {$1$}}
            }
            child {node {$)$}}
        }
        child {node {$AND$}}
        child {node {$E$}
            child {node {$1$}}
        }
        child {node {$)$}};
    \end{tikzpicture}
    \caption{Esempio di albero di derivazione}
\end{figure}

\noindent In tal merito, data una grammatica $G$ libera da contesto, vale $S\xRightarrow{G}_* \alpha$ se e solo se esiste un albero di derivazione con
radice $S$ e sequenza di foglie $\alpha$. Siccome un albero di derivazione rappresenta un certo insieme di possibili derivazioni di una stessa stringa,
ciò implica che possano esistere più alberi per la stessa parola e quando questi condividono la radice $S$ diciamo che la grammatica CF è
\textbf{ambigua}. Eventuali linguaggi che possono essere generati solamente da tali grammatiche si chiamano \textbf{inerentemente ambigui}.

%

\section{Semplificazione delle grammatiche CF}
Lavorare con grammatiche complesse può risultare oneroso, per questa ragione sono stati creati alcuni modi per restringere la forma delle produzioni a
loro permesse senza alterare la classe dei linguaggi che riconoscono. In particolare, queste trasformazioni hanno lo scopo di raggiungere \textbf{forme
normali}, le quali descrivono le produzioni con vincoli di forma su $\alpha$. Ne esistono di due tipi: \begin{itemize}
    \item \textbf{Chomsky}: $A\to BC, A\to a$, dove $B,C\in V, a\in T$.
    \item \textbf{Greibach}: $A\to a\alpha$, dove $a\in T, \alpha\in (T\cup V)^*$.
\end{itemize}

\noindent Ciò implica che per ogni grammatica CF tutte le produzioni possono essere trasformate in forma normale, siccome ogni linguaggio CF può essere
generato da una grammatica in forma normale. Per arrivarci, è necessario eseguire tre operazioni di \textbf{semplificazione}.

\paragraph{Eliminazione dei simboli inutili:} Data una grammatica $G = \left<V, T, P, S\right>$ diciamo che un generico simbolo $x$ terminale o non
terminale è utile quando esiste almeno una derivazione dal simbolo iniziale $S$ a una sequenza di simboli terminali $\omega$ che ne fa uso. Dunque: \[
    S\xRightarrow{}_* \alpha x\beta \xRightarrow{}_* \omega;\quad \text{Dove } \alpha, \beta \in (V\cup T)^*, \omega\in T^*
\]

\noindent Questa operazione si divide in due parti: in primo luogo si fa una rimozione dei simboli non terminali che non raggiungono simboli terminali,
alla quale segue la rimozione dei simboli non terminali non raggiunti da quello iniziale $S$. Supponiamo quindi una grammatica
$G = \left<\{A,B,C,D,E,S\}, \{0,1,2,3\}, P, S\right>$, dove: \[
    V = \{A,B,C,D,E,S\}\quad T = \{0,1,2,3\}
\]

\noindent Le cui produzioni $P$ sono: \begin{itemize}
    \item $S \to \epsilon|A|B|D$
    \item $A \to 0|0A0|B|E$
    \item $B \to 1|1B1|A|E$
    \item $C \to A|B|2$
    \item $D \to \epsilon|D$
    \item $E \to \epsilon$
\end{itemize}

\noindent L'idea è di partire dai simboli terminali e verificare all'indietro nelle produzioni la presenza dei vari simboli non terminali. Quelli che
sono indicati nelle produzioni ma non sono ottenuti con questo ragionamento sono inutili. Dunque, costruiamo l'insieme di simboli non terminali
$\Gamma^i$ iterativamente, i cui elementi raggiungeranno almeno un simbolo terminale in un numero $i$ di passi. Notiamo che le produzioni che contengono
simboli terminali sono $A, B, C$, quindi aggiungiamole a $\Gamma^1$; è possibile vedere anche che $S$ genera $A$ e $B$, quindi va aggiunto a $\Gamma^2$.
Non sono presenti altri simboli utili per la prima fase; risulta quindi: \[
    \Gamma^0 = \emptyset\quad \Gamma^1 = \{A, B, C\}\quad \Gamma^2 = \{A, B, C, S\}
\]

\noindent Da qui otteniamo la nuova grammatica $G' = \left<V\cap \Gamma^2, T, P', S\right>$ con nuove produzioni $P' = P / \{D\to \epsilon | D,
E\to \epsilon\}$. Procediamo ora con la rimozione dei simboli non terminali che non sono raggiunti dal simbolo iniziale $S$; abbiamo ottenuto le nuove
produzioni: \begin{itemize}
    \item $S \to \epsilon|A|B$
    \item $A \to 0|0A0|B$
    \item $B \to 1|1B1|A|\epsilon$
    \item $C \to A|B|2$
\end{itemize}

\noindent Sempre iterativamente, creiamo l'insieme $\Gamma^i$ dei simboli raggiunti da $S$ in $i$ passi: \[
    \Gamma^0 = \{S\}\quad \Gamma^1 = \{A, B, S\}\quad \Gamma^2 = \{A, B, S, 0, 1\}
\]

\noindent Otteniamo un'ulteriore grammatica $G'' = \left<V'\cap \Gamma^2, T' \cap \Gamma^2, P'', S \right>$, dove $V'$ è l'insieme dei non terminali da
cui è stato rimosso $C$, in quanto non generato da $S$, e $T'$ è l'insieme dei terminali da cui sono stati rimossi $2,3$. Ne consegue che le produzioni
$P''$ sono: \begin{itemize}
    \item $S \to \epsilon|A|B$
    \item $A \to 0|0A0|B$
    \item $B \to 1|1B1|A|\epsilon$
\end{itemize}

% TODO Riprendi da lez9

\paragraph{Eliminazione delle produzioni unitarie:}

\paragraph{Eliminazione delle $\epsilon$-produzioni:}

%

\section{Automi a pila}

%

\section{Classificazione di Chomsky}